% Πανεπιστήμιο Θεσσαλίας
% Υπόδειγμα Εργασίας Σχεδιασμού Ενεργειακών Συστημάτων
% Τμήμα Μηχανολόγων Μηχανικών
%------------------------------------------------
\documentclass[a4paper,oneside,12pt]{book}
 
\usepackage[LGR, T1]{fontenc}
\usepackage[english,greek]{babel}
%\usepackage{alphabeta}
\usepackage{newtxtext}
\usepackage{substitutefont}
\substitutefont{LGR}{\rmdefault}{txr}

\usepackage{amsbsy, amsmath, amsfonts}
\usepackage{siunitx}
\usepackage{graphicx}

% text sub- and superscripts
\newcommand{\spr}[1]{\ensuremath{^\textrm{#1}}}
\newcommand{\sub}[1]{\ensuremath{_\textrm{#1}}}

% better tables
\usepackage{multirow}
\usepackage{booktabs}

% links
%\usepackage{color}
\usepackage[]{hyperref}
\usepackage{enumitem} % better lists
%βιβλιογραφία
\usepackage[backend=biber,style=numeric,sorting=none,natbib=true, babel=other,bibencoding=auto, language=autobib,giveninits=true]{biblatex}
\addbibresource{myrefs.bib} %file of biblatex
\newcommand{\sw}{\selectlanguage{english}}
\newcommand{\sq}{\selectlanguage{greek}}
\newcommand{\eng}[1]{\latintext#1\greektext}
\newcommand{\gre}[1]{\greektext#1\latintext}
% figures and tables
\usepackage{caption}
\captionsetup[figure]{labelfont=it,textfont=it}
\captionsetup[table]{labelfont=it,textfont=it}
\usepackage{subfig}

\usepackage[strict]{changepage}

% Define margins
\usepackage[left=2.5cm,right=2.5cm,top=2.5cm,bottom=2.5cm,includehead,includefoot,headheight=16pt]{geometry}

% Page headers
\usepackage{fancyhdr}
\fancyhf{}
\renewcommand{\headrulewidth}{0pt} 
\fancyhead[L]{\nouppercase{\leftmark}}
\cfoot{\thepage}
\pagestyle{fancy}


% Line spacing
\usepackage{setspace}
\setstretch{1.5}

% Defines uth-specific title page
\usepackage{mieUTHtitle}


%Σωστός συλλαβισμός λέξεων. Εάν ο συλλαβισμός μια λέξης δεν είναι σωστός και το LaTex βάζει παύλα σε λάθος σημείο στην αλλαγή γραμμής τότε δίνουμε εδώ τον σωστό συλλαβισμό.
\babelhyphenation[greek]{Α-ντι-κεί-με-νο}
%-------------------------------------------
\begin{document}

% Συμπληρώστε τα στοιχεία σας
\author{Όνομα Επώνυμο}
\title{Τίτλος}

\maketitle

\frontmatter

\chapter*{Περίληψη}
Η περίληψη θα πρέπει να είναι μέχρι 300 λέξεις. Θα πρέπει να συνοψίζει το αντικείμενο της εργασίας, τις μεθόδους που χρησιμοποιήθηκαν, τα βασικά σημεία της μελέτης, τα ευρήματα και τα συμπεράσματα.

\tableofcontents

%-------------------------------------------------
\mainmatter

\chapter{Εισαγωγή}

\section{Μορφοποίηση Πινάκων, Σχημάτων και Εξισώσεων}
Κάθε αντικείμενο πρέπει να είναι αριθμημένο και να έχει τίτλο/περιγραφή  (λεζάντα). Οι λεζάντες των εικόνων πρέπει να εμφανίζονται κάτω από αυτές, για παράδειγμα δες Σχήμα \ref{fig:wave}. Οι λεζάντες των πινάκων πρέπει να εμφανίζονται πάνω από αυτούς, για παράδειγμα δες τον Πίνακα \ref{tab:my_label}. Η αρίθμηση των εξισώσεων εμφανίζεται δεξιά τους (ή από κάτω εάν ο χώρος δεν είναι επαρκής), δες Εξ.\ \ref{eq:y}.

Μία αποσβένουσα αρμονική δίνεται ως,
\begin{equation}\label{eq:y}
    y(t)=a e^{-b\omega t} \sin(\omega t + \phi)
\end{equation}
όπου $a$ είναι εύρος του ημιτονοειδούς κύματος, $\omega$ η συχνότητα,  $\phi$ η μετατόπιση φάσης και ο εκθετικός συντελεστής $b>0$. Ένα παράδειγμα δίνεται στο Σχήμα \ref{fig:wave}. 

\begin{figure}[thb]
	\centering
	\includegraphics[width=0.9\textwidth]{figure_wave.eps}
	\caption{Εκθετική απόσβεση ημιτονοειδούς κύματος από την Εξ.\ \ref{eq:y}}\label{fig:wave}
\end{figure}

Άλλα παραδείγματα εξισώσεων
\begin{equation}\label{eq:opt}
\min\limits_{u,x} \;\phi\left(x(t_f)\right) \;+\int_{t_0}^{t_f} f_0\left(x(t),u(t)\right)\,dt
\end{equation}
εξαρτόμενη από
\begin{align}
\dot{x}(t)&=f\left(x(t),u(t)\right)\\
c\left(x(t),u(t)\right) &\leq 0,\;\; t\in[t_0,t_f]\\
\omega\left(x(t_0),x(t_f)\right) &= 0 
\end{align}
όπου $f(x(t), u(t))$ ..., η συνάρτηση $\omega(x(t_0),x(t_f))$ δίνει ...και  $c(x(t),u(t))$ είναι περιορισμοί.

Πίνακες, σχήματα και εξισώσεις πρέπει να έχουν ξεχωριστή αρίθμηση. Κάθε αντικείμενο πρέπει να εμφανίζεται με τη σωστή σειρά και να έχει αναφερθεί στο κυρίως κείμενο πριν εμφανιστεί. Εάν χρησιμοποιείσεται σχήματα ή πίνακες από άλλες πηγές (βιβλία, δημοσιεύσεις, ιστοσελίδες, κτλ.) τότε πρέπει να αναφέρεται στη λεζάντα η πηγή όπως στους Πίνακες \ref{tab:my_label}--\ref{tab:my_label2}. Εάν έχετα συνθέσει ένα σχήμα ή πίνακα από διάφορες πηγές τότε πρέπει να αναφέρονται είτε μέσα στο σχήμα/πίνακα είτε στη λεζάντα του.

Τα αντικείμενα πρέπει να είναι στιχοισμένα στο κέντρο. Τυχών κείμενο μέσα στα σχήματα πρέπει να είναι ευανάγνωστο.Δεν υπάρχει περιορισμός στο των σχημάτων, καλύτερα πάντως να δείξετε μία μεγάλη και ευκρινή εικόνα, παρά μία μικρότερη και δυσανάγνωστη.


\begin{table}[htb]
    \centering
    \begin{tabular}{l l l p{5cm}}
    \toprule
    Μέρα & Ελ. Θερμ. & Μεγ. Θερμ. & Σύνοψη \\ \midrule
    Δευτέρα & \SI{11}{\celsius} & \SI{22}{\celsius} & Μία καθαρή μέρα με αρκετή ηλιοφάνεια. Όμως οι ισχυροί άνεμοι θα ρίξουν την θερμοκρασία. \\ \midrule
    Τρίτη & \SI{9}{\celsius} & \SI{19}{\celsius}  & Συννεφιά με βροχή κατά μήκος των βόρειων περιοχών. \\ \midrule
    Τετάρτη & \SI{10}{\celsius}  & \SI{21}{\celsius}  & Οι συνθήκες θα βελτιωθούν από νωρίς το απόγευμα. \\
    \bottomrule
    \end{tabular}
    \caption{Λεζάντα πίνακα \cite{website:table}}
    \label{tab:my_label}
\end{table}

\begin{table}[htb]
    \centering
\begin{tabular}{llr}  
\toprule
\multicolumn{2}{c}{Αντικείμενο} \\
\cmidrule(r){1-2}
Ζώο    & Περιγραφή & Τιμή (\$) \\
\midrule
Γκνατ      & το γραμμάριο    & 13.65      \\
          &    τεμάχιο     & 0.01       \\
Γκνου       & γεμιστό     & 92.50      \\
Εμού     & γεμιστό     & 33.33      \\
Αρμαντίλλο & κατεψυγμένο      & 8.99       \\
\bottomrule
\end{tabular}
    \caption{Λεζάντα πίνακα \cite{website:table}}
    \label{tab:my_label2}
\end{table}

\textlatin{Latin text can also be added to 
the document.}


\section{Αναφορές - Βιβλιογραφία}
Είναι σημαντικό να δώσετε προσοχή στη σωστή αναφορά της δουλειάς άλλων (συμπεριλαμβανομένων των σχημάτων), τόσο για την αποφυγή της λογοκλοπής όσο και για να επιτρέψετε σε άλλους να βρουν το πηγαίο κείμενο \cite{blackholes}, π.χ., σύνδεσμοι του διαδυκτίου θα πρέπει να περιέχουν την πλήρη διεύθυνση \textlatin{URL} καθώς και την ημερομηνία επίσκεψης της ιστοσελίδας. Το λογισμικό \textlatin{TurnItIn} θα χρησιμοποιηθεί για τον έλεγχο τυχών λογοκλοπής.

Στη βιβλιογραφεία μπορούν να αναφερθούν διαφορετικές πηγές, συμπεριλαμβανομένων βιβλίων \cite{book:goossens93} και άρθρων \cite{art:AbedonHymanThomas2003,art:greenwade93,proc:Abedon1994} ή ακόμη και πηγών στα ελληνικά \cite{KRitos2021}.



%--------------------------------------------
\chapter{Βιβλιογραφική Ανασκόπηση}


%--------------------------------------------
\chapter{Σχεδιασμός \& Θερμοδυναμική Ανάλυση }
\section{Ενότητα 1}

\section{Ενότητα 2}
\subsection{Υποενότητα}

\section{Ενότητα 3}

%--------------------------------------------
\chapter{Οικονομική Ανάλυση}


%--------------------------------------------
\chapter{Επίλογος}
Αυτό το κεφάλαιο συνοψίζει τις πληροφορίες που παρουσιάστηκαν στην τεχνική έκθεση. Καμία καινούργια πληροφορία δεν πρέπει να παρουσιάζεται εδώ. Συχνά μπορεί να περιλαμβάνει μία ενότητα με περαιτέρω ή μελλοντική εργασία, όπου παρουσιάζεται τι θα ήταν χρήσιμο να γίνει στο μέλλον.

%----------------------------------------------
% Inserts references/bibliography, using myrefs.bib as a database but only includes those entries referenced in the text above
\selectlanguage{english}
\printbibliography[title={\textgreek{Βιβλιογραφία}}]
\end{document}